% !TEX root = ../Tesis.tex
\chapter{Discusión de resultados}
\section{Condiciones atmosféricas y calidad del aire}

Durante la campaña del 20 de febrero al 5 de marzo de 2026, las concentraciones de los contaminantes 
criterio en las estaciones de monitoreo se mantuvieron, en general, por debajo de los valores de referencia
de la normatividad mexicana (Tablas 3.5 y 3.6). El ozono alcanzó un máximo diario de 45~ppb 
(Cruz Roja, 23 de febrero); el SO$_2$ se mantuvo entre aproximadamente 10 y 20~ppb en DIF y Nativitas, 
frente al límite de 0.040~ppm para 24~h; y el PM$_{2.5}$ llegó a cerca de 27~$\mu$g\,m$^{-3}$ 
(DIF, 20 de febrero), por debajo de los 41~$\mu$g\,m$^{-3}$ del primer nivel de la norma. La excepción 
fue el PM$_{10}$, que alcanzó cerca de 100~$\mu$g\,m$^{-3}$ en Nativitas el 22 de febrero, por encima del 
valor de 70~$\mu$g\,m$^{-3}$. Sin embargo, la norma se evalúa con el percentil 99 anual de los promedios 
de 24~h, por lo que un periodo de dos semanas no permite dictaminar cumplimiento: estos valores caracterizan
las condiciones del muestreo, no el estado normativo de la ciudad. 

Las diferencias entre estaciones son sistemáticas. Nativitas presentó las concentraciones más 
altas de SO$_2$, PM$_{10}$ y CO; Cruz Roja, las de NO$_2$; y DIF, las de PM$_{2.5}$ 
(Nativitas no reportó PM$_{2.5}$). A escala estatal, el inventario de emisiones de 2021 atribuye 
casi todo el SO$_2$ a fuentes fijas (7\,331.79 de 7\,426.74~t/año) y la mayor parte del CO a 
fuentes móviles (23\,411.99 de 26\,974.73~t/año) (Tabla 3.1). El patrón de Nativitas sugiere 
entonces una combinación de influencia industrial y vehicular, mientras que el de Cruz Roja podría 
responder a una mayor influencia del tránsito. Lo relevante para el objetivo 
isotópico es que las retrotrayectorias de los cuatro receptores son prácticamente coincidentes 
(Sección 5.5), de modo que estas diferencias reflejan fuentes locales y no un transporte regional distinto.

Los ciclos diurnos muestran un patrón común. El NO$_2$ y el CO presentan dos máximos, hacia las 8 y las 
20~h, y el SO$_2$, el PM$_{10}$ y el PM$_{2.5}$ alcanzan sus máximos entre las 20 y las 21~h en Nativitas 
y DIF. Esta estructura es compatible con emisiones vehiculares en horas pico que se acumulan bajo una capa 
límite poco profunda, y sitúa el máximo vespertino inmediatamente antes de la ventana de muestreo 
(22:00--02:00~h). Como la capa límite nocturna es estable y limita la mezcla vertical (Sección 2.2), 
los contaminantes emitidos al final de la tarde tienden a permanecer cerca de la superficie durante la 
noche, lo que hace de esa ventana un momento favorable para detectar vapor derivado de la combustión. 
El ozono, con máximo a las 17~h y mínimo a las 7~h, sigue el ciclo fotoquímico esperado.

En cuanto a la meteorología, las rosas de viento de las estaciones muestran vientos predominantes 
del sureste de 2.1 a 3.6~m\,s$^{-1}$, mientras que la rosa del periodo de muestreo señala vientos del 
noreste, con velocidades de 2.1 a 5.7~m\,s$^{-1}$. Estas velocidades moderadas implican una ventilación 
mayor que la de los episodios de estancamiento de las ciudades de Estados Unidos y China en las que se 
basan los trabajos de referencia, caracterizadas por fuertes inversiones térmicas invernales (Sección 1.2), 
lo que es relevante para interpretar la magnitud de la señal isotópica (Sección 6.2). 


\section{Firma isotópica del vapor y contribución del vapor de combustión}

Las muestras de vapor de agua colectadas en Salamanca presentan valores de $\delta^{18}$O entre 
$-6$ y $-12$\,\textperthousand\ y de $\delta^{2}$H entre $-75$ y $-120$\,\textperthousand, y se
alinean en torno a la isolínea de mezcla de 20\,\% del diagrama $\delta^{2}$H--$\delta^{18}$O (Figura 5.17),
sin que ninguna supere el 30\,\%. El modelo de mezcla binaria detecta, por tanto, una contribución de vapor 
derivado de la combustión (VDC), pero minoritaria, en la dirección esperada desde la línea meteórica hacia 
el extremo de combustión (Gorski et al., 2015; Xing et al., 2020). Un rasgo informativo es que los puntos 
se distribuyen a lo largo de una misma isolínea y no a través de varias: la variabilidad de $\delta^{2}$H 
y $\delta^{18}$O entre muestras parece reflejar sobre todo cambios en el vapor de fondo, más empobrecido o 
más enriquecido según la masa de aire, y no cambios en la fracción de VDC. 

Que la fracción de VDC sea moderada es razonable para este contexto. Los trabajos de referencia se 
realizaron en ciudades con inversiones térmicas intensas y episodios prolongados de estancamiento, 
condiciones que concentran las emisiones en la capa límite y favorecen que la señal de combustión destaque 
sobre un vapor ambiental de baja humedad (Yang y Yoshimura, 2024). En Salamanca, durante la campaña, los
vientos fueron moderados, por lo que una fracción cercana al 20\,\% no resulta una anomalía. A pesar de eso, 
es una cantidad apreciable de vapor de origen antrópico para un municipio de unos 273 mil habitantes, y es 
coherente con la concentración de fuentes de combustión en la zona: la central termoeléctrica, la refinería 
y las industrias de cerámicos y automotriz registradas en el RETC (Tablas 3.2 y 3.3). Debido a la hora de 
muestreo, otras fuentes como el tránsito vehicular, la quema de esquilmos y las ladrilleras (Capítulo 3) 
pueden considerarse despreciables. El modelo no distingue entre estas fuentes: el extremo de combustión 
es un único punto teórico, mientras que la composición isotópica del VDC varía con el tipo de combustible.

La fracción estimada depende de dos supuestos que conviene discutir. El primero es el extremo de 
combustión, fijado en $\delta^{18}$O $\approx +9$\,\textperthousand\ y
$\delta^{2}$H $\approx -135$\,\textperthousand, valor teórico y no medido en Salamanca. Como el VDC 
hereda su $\delta^{18}$O del oxígeno atmosférico ($\approx +23.9$\,\textperthousand) y su $\delta^{2}$H del 
combustible (Sección 2.4.1), el extremo variará con el combustible y con el fraccionamiento durante la 
combustión, y cualquier cambio en él modifica el porcentaje de mezcla. El segundo es el vapor de fondo, 
definido por la línea LM-SDR ($\delta^{2}$H $= 7.289\,\delta^{18}$O $+ 4.4899$), ajustada con datos de 
San Juan del Río y no de Salamanca. Su pendiente es casi igual a la de la línea meteórica local de la 
Cuenca Alta del Río Laja ($7.349$; Ramírez-González et al., 2024), pero la ordenada difiere en unos 
3.8\,\textperthousand, lo que desplaza el punto de fondo asignado a cada muestra. 

\section{Modelos de dispersión con FLEXPART}

Las retrotrayectorias promedio nocturnas (22:00--02:00~h) de los cuatro receptores son similares: 
descienden desde 3.0--3.5~km hasta 1.5--2.0~km en los últimos $\sim$130~km y se proyectan hasta poco más 
de 400~km del receptor en dirección noreste (Figura 5.18). Por tratarse de retrotrayectorias, esto indica 
que el aire llegó a Salamanca desde esa dirección, en concordancia con la rosa de viento del periodo de
muestreo (Figura 5.16). 

Este origen común tiene dos implicaciones. La primera es que el vapor de fondo que llega a los cuatro 
puntos es esencialmente el mismo, de modo que las diferencias entre estaciones (Sección 6.1) se atribuyen 
a emisiones locales. La segunda es que el flujo descendente implica subsidencia, que tiende a estabilizar 
los niveles bajos y a aportar aire probablemente más seco, condiciones que limitan la dilución del VDC 
emitido en superficie y favorecen su detección.

Los \textit{footprints} de CO refuerzan esta lectura local. En los cuatro receptores, la sensibilidad 
más alta ($>0.3$~u.a.) se concentra en una franja continua contigua al receptor, con máximos ($>0.5$~u.a.) 
en las celdas inmediatas, y decae a menos de 0.1~u.a. en la mayor parte del dominio (Figuras 5.20 y 5.21). 
El CO medido en cada punto procede, por tanto, sobre todo de emisiones próximas, dentro del área urbana, y 
no de fuentes distantes. Si el vapor de combustión se emite junto con el CO, esto sugiere que la señal 
isotópica refleja principalmente combustión local. La inferencia es indirecta: el CO es un trazador de 
combustión, pero el vapor de agua tiene además fuentes naturales y no se emite en proporción fija con 
el CO, de modo que el \textit{footprint} describe la región fuente del CO y solo es un indicador 
aproximado de la del VDC. La zona de alta sensibilidad es más extensa en la ventana nocturna que en el 
promedio de 24~h. 





